\documentclass[11pt,a4paper]{article}
\usepackage[T1]{fontenc}
\usepackage{lmodern,amsmath,amssymb,amsthm,mathtools,microtype}
\usepackage[margin=27mm]{geometry}
\usepackage[colorlinks=true,linkcolor=blue,citecolor=blue,urlcolor=blue]{hyperref}
\setlength{\emergencystretch}{2em}
\newtheorem{theorem}{Theorem}[section]
\newtheorem{lemma}[theorem]{Lemma}
\newtheorem{corollary}[theorem]{Corollary}
\newcommand{\Q}{\mathbb Q}\newcommand{\Z}{\mathbb Z}
\newcommand{\F}{\mathbb F}
\title{Opposite signs after duplicating elliptic reference points}
\author{}\date{16 November 2026}
\begin{document}\maketitle
\begin{abstract}
An explicit progression of global elliptic reference points has opposite
formal-parameter translation signs after duplication. The points lie in
the complete reference domain used by the quartic constructions, whose
indices may be half-integral. The relevant order-four residue point is
itself a finite-field double, although it is outside the reduction of
the subgroup generated by the distinguished doubled point. Exact order
certificates establish the progression at a permissible soundness contact
prime. The result concerns reference branches; it does not assert a
solution of the full detector equations.
\end{abstract}

\section{A global reference progression}

Put
\[
 K=\Q(\rho),\qquad \rho^2=2,\qquad D=11-6\rho,
 \qquad E:y^2=x^3-Dx,
\]
\[
 R_0=(1+\rho,2\rho),\qquad
 P=2R_0=\left(\frac{57-28\rho}{8},
                  \frac{-660+417\rho}{32}\right),
 \qquad T_0=(0,0).
\]
For a root $r$ of $X^2-2$ modulo a good rational prime $p$, write
$\operatorname{red}_{p,r}$ for reduction at $(p,\rho-r)$.
The complete reference chart used in \cite{ABC79,ABC78,Chart}
contains every point in
\begin{equation}\label{eq:domain}
                    E(K)\setminus\{O,T_0,P,P+T_0\}.
\end{equation}
In particular its domain is not restricted to $\Z P$.

\begin{samepage}
\begin{theorem}\label{thm:progression}
Let
\[
 p=76537,\qquad r_1=26935,\qquad r_2=49602,
\]
and, for every integer $j\ge0$, set
\[
 m_j=10435875+61021500j,\qquad H_j=[m_j]R_0.
\]
The points $H_j$ are distinct and belong to~\eqref{eq:domain}. They satisfy
\[
 \operatorname{red}_{p,r_1}(H_j)=O,
 \qquad
 \operatorname{red}_{p,r_2}(H_j-P)=(35686,1608)=:R,
\]
where $[2]R=T_0$ and $R$ has exact order four. Consequently
\[
 \operatorname{red}_{p,r_1}([2]H_j)=O,
 \qquad
 \operatorname{red}_{p,r_2}([2](H_j-P))=T_0.
\]
Both duplicated characteristic-zero points are non-torsion.
\end{theorem}
\end{samepage}

The index relative to $P$ is $m_j/2$, a half-integer with odd numerator.
These are permitted reference indices in the quartic sources.
Under their retained assertion $2E(K)=\Z P$, none of the $H_j$
belongs to $2E(K)$: an equality $[m_j]R_0=[n]P$ would contradict
non-torsion of $R_0$ and oddness of $m_j$.

\section{Exact finite certificates}

The two reductions have the following data:
\[
\begin{array}{c|r|r|r|r}
 \rho&D&R_0&\operatorname{ord}(R_0)&\operatorname{ord}(P)\\\hline
 26935&68012&(26936,53870)&6375&6375\\
 49602&8547 &(49603,22667)&9572&4786.
\end{array}
\]
Trial division through $\lfloor\sqrt{76537}\rfloor=276$ proves
primality. The two displayed roots square to $2$ and are opposite
modulo $p$. Both values of $D$ are nonzero, so both reductions are good.

Write $R_{0,1},R_{0,2}$ for the reduced points. Direct addition gives
$[6375]R_{0,1}=O$ and the nonzero proper prime-quotient multiples
\[
\begin{array}{c|r}
 \text{multiplier}&\text{point}\\\hline
 6375/3 &(7526,56152)\\
 6375/5 &(70736,15850)\\
 6375/17&(21124,73399).
\end{array}
\]
Thus the first order is exact. In the second reduction,
\[
 [9572]R_{0,2}=O,\qquad
 [4786]R_{0,2}=T_0,\qquad
 [4]R_{0,2}=(46317,70872).
\]
Since $9572=4\cdot2393$ and $2393$ is prime, these identities prove
the second order. They also give
\begin{equation}\label{eq:quarter}
 [2393]R_{0,2}=(35686,1608),\qquad
 [2](35686,1608)=T_0.
\end{equation}
Since $P=2R_0$, its order is
$\operatorname{ord}(R_0)/\gcd(2,\operatorname{ord}(R_0))$,
giving the stated values.

The point $P$ is non-torsion by the elementary good-reduction argument
of \cite[Lemma~2.1]{Mixed}: its first displayed reduction has order
divisible by $3$, whereas its reduction at $(17,\rho-6)$ has order
$5$. A nonzero $3$-primary torsion point cannot lie in the formal
kernel at a good prime of residue characteristic $17$.
This proof uses no group-rank assertion.

\begin{proof}[Proof of Theorem~\ref{thm:progression}]
The two orders of $R_0$ are coprime and have product $61021500$.
The initial coefficient satisfies
\[
 10435875=6375\cdot1637,
 \qquad10435875-2=9572\cdot1090+2393.
\]
The same congruences hold for every $m_j$. Reduction is a group
homomorphism, so the order certificates and~\eqref{eq:quarter} give
the asserted reductions. The coefficients $m_j$ and $m_j-2$ are
nonzero, and $R_0$ is non-torsion. Thus $H_j,H_j-P$, and their
doubles are non-torsion. In particular the excluded points
in~\eqref{eq:domain} cannot occur. The distinct positive $m_j$
give distinct global points.
\end{proof}

The finite-field double subgroup is a separate condition. At the second
place, exact addition gives
\begin{equation}\label{eq:finitehalf}
                   [2](24152,13484)=(35686,1608)=R.
\end{equation}
Hence $R\in2E(\F_p)$. But it does not belong to the subgroup generated
by the reduced $P$, whose order $4786$ is not divisible by four.
Under $2E(K)=\Z P$, this proves that the reduction of the global
double subgroup is a proper subgroup of $2E(\F_p)$ at this place.

\section{The signs after duplication}

Use the rational parameter $z=-x/y$, with its regular extensions at
$O$ and $T_0$. Translation by $T_0$ gives
\[
 (x,y)\longmapsto\left(-\frac D x,\frac{Dy}{x^2}\right),
 \qquad z(Q+T_0)=-z(Q).
\]
These are identities of rational functions on the smooth curve.

The point $R$ in the second reduction has a unique order-four lift
$\widetilde R$ with $[2]\widetilde R=T_0$ over the corresponding
completion of $K$. For a direct construction, lift its abscissa by
$x^2+D=0$ and then its chosen ordinate by $y^2=x^3-Dx$.
Both derivatives are units at the displayed nonzero residue coordinates,
so simple-root Hensel lifting applies. The doubling formula gives
$[2]\widetilde R=T_0$.

\begin{corollary}\label{cor:signs}
For every $j$, the point $[2]H_j$ is a nonidentity formal point at
the first completion. At the second completion put
$Q_j=H_j-P-\widetilde R$. Then $Q_j$ is a nonidentity formal point and
\[
 z([2](H_j-P))=z(T_0+[2]Q_j)=-z([2]Q_j).
\]
Thus the two reference neighborhoods give opposite translation signs
even after duplication.
\end{corollary}
\begin{proof}
The reduction statements place the indicated points in the formal
groups. If $Q_j=O$, the global point $H_j-P$ would be torsion, which
was excluded. Its double is also nonzero. The identity follows from
$[2]\widetilde R=T_0$ and the displayed translation formula.
\end{proof}

\section{Contact-prime and detector scope}

The soundness argument in \cite{ABC79,ABC78} consumes every prime
$p>76176$ at which a contact has positive odd valuation. It does not
restrict those primes to the Gaussian ray classes used to construct
witnesses in completeness. Our prime is a concrete such contact prime.
For
\[
 \Theta=\{-5,-2,-1,0,2,3,4,7\},\qquad
 F(t)=\prod_{\theta\in\Theta}(t-\theta),\qquad
 w(t)=\frac{F(t)}{(t-1)^6},
\]
one has $F'(0)=1680\not\equiv0\pmod p$. The rational argument
$t=p$ therefore satisfies $v_p(w(t))=1$. All other contact factors
and $t-1$ are units. The prime is split in $K$ and has good reduction.

It is not in the stronger selected Gaussian ray class: the fixed
residue $1+3i$ modulo seven in \cite[Section~2]{Dependencies} has norm
$3$ modulo seven, whereas $76537\equiv6\pmod7$. Those restrictions
belong to completeness, so they do not exclude the present branch
from an unconditional soundness-prime analysis.

The progression imposes actual global reference points, their fixed
translation by $P$, and the complete reference chart. It does not
impose the numerator points, ratio comparisons, second detector argument,
height equations, or pole witnesses. In particular the result does not
contradict the defining-polynomial soundness theorems, whose full local
tests impose additional conditions. It excludes automatic agreement of
the duplicated reference signs from the point domain and admissible
contact primes alone.

\section{Reproducible verification}

Run
\begin{center}\small
\texttt{python negative\_results/computation/duplicated\_sign\_v44/check.py}.
\end{center}
The standard-library checker reconstructs the points and all exact
order certificates using separate affine and Jacobian implementations.
It checks trial-division primality, the progression congruences,
the finite-field half in~\eqref{eq:finitehalf}, and the exact rational
contact valuation. It binds the proof and its source files by SHA-256.
The formal lifting and infinite progression arguments are the proofs
above, rather than extrapolations from finite examples.

\begin{thebibliography}{9}
\raggedright
\bibitem{ABC79}\emph{A 79-variable quartic under rational $abc$},
project source, version27,
\texttt{papers/quartic\_existence/}\allowbreak
\texttt{quartic\_abc79\_20261114\_v27.tex}.
\bibitem{ABC78}\emph{A quartic in 78 variables from a mixed-power height estimate},
project source, version34,
\texttt{papers/quartic\_existence/}\allowbreak
\texttt{quartic\_abc78\_20261115\_v34.tex}.
\bibitem{Chart}\emph{A three-coordinate quadratic chart for an elliptic point and its translate},
project source, version40,
\texttt{papers/elliptic\_charts/}\allowbreak
\texttt{three\_coordinate\_reference\_20261115\_v40.tex}.
\bibitem{Mixed}\emph{Mixed torsion signs for doubled elliptic reference points},
project source, version42,
\texttt{negative\_results/papers/elliptic\_charts/}\allowbreak
\texttt{mixed\_reference\_signs\_20261116\_v42.tex}.
\bibitem{Dependencies}\emph{Arithmetic dependencies of the construction with 29 variables},
project source, version8,
\texttt{papers/arithmetic\_dependencies/}\allowbreak
\texttt{arithmetic\_dependencies\_20261108\_v8.tex}.
\end{thebibliography}
\end{document}
